\documentclass[11pt]{article}

\usepackage[margin=1in]{geometry}
\usepackage{amsmath,amssymb,amsthm}
\usepackage{booktabs}
\usepackage{microtype}
\usepackage[hidelinks]{hyperref} \hypersetup{pdftitle={The Exact Matching Spectrum of a Three-Job Wang--Sitters Tight Family},pdfauthor={Xiaojian Zeng}}

\newtheorem{theorem}{Theorem}

\title{The Exact Matching Spectrum of a Three-Job\\Wang--Sitters Tight Family}
\author{Xiaojian Zeng\\
\small PBC School of Finance, Tsinghua University, 100190\\
\small \texttt{xiaojianzeng@foxmail.com}}
\date{}

\begin{document}
\maketitle

\begin{abstract}
For the three-job, three-machine restricted-assignment family used to approach
the $11/6$ bound for Wang--Sitters rounding, we determine every valid slot
matching for every real parameter $0<\delta<1/6$.  The half-open slot graph is
constant throughout this interval and has exactly eight job-saturating
matchings.  Their makespans, counted with matching multiplicity, form the
multiset
\[
 \left\{\left(\frac{11}{6}-\delta\right)^1,
 \left(\frac43-\delta\right)^2,1^5\right\}.
\]
Thus the minimum makespan is $1$ and the maximum is
$11/6-\delta$.  This is a symbolic census of one fixed fractional family; it
does not change the approximation guarantee for arbitrary instances.
\end{abstract}

\section{Context and statement}

The matching-based rounding of Shmoys and Tardos
\cite{ShmoysTardos1993} is a basic tool for assignment problems.  Wang and
Sitters used this rounding in a simple $11/6$-approximation for graph balancing
\cite{WangSitters2016}.  Shavit later isolated the nondeterminism in the
fractional solution and in the valid slot matching \cite{Shavit2026}.  His
Theorem~2 gives the parametric family studied below and one valid matching of
makespan $11/6-\delta$ for every $0<\delta<1/6$.  His Measurement~4 reports,
at $\delta\in\{10^{-2},10^{-3},10^{-4}\}$, eight valid matchings with minimum
$1$ and maximum $11/6-\delta$.  The result here upgrades that finite
measurement to a proof uniform on the full open interval.

Fix $0<\delta<1/6$, machines $\{0,1,2\}$, and the three jobs
\[
\begin{array}{c|c|c}
 j & p_j & \{(i,x_{ji}):x_{ji}>0\}\\ \hline
 q&1&\{(0,\frac13+\delta),(1,\frac23-\delta)\}\\
 h&\frac12&\{(0,\frac23),(2,\frac13)\}\\
 r&\frac13-\delta&\{(0,1)\}.
\end{array}
\]
The threshold is $T=1$ and the Step-2 threshold is $\beta=2/3$.
A job is big when its size is strictly greater than $1/2$, so only $q$ is
big.  Both shares of $q$ are strictly below $2/3$, and hence Step~2 assigns no
job.

On each machine, order the surviving fractional jobs by nonincreasing size and
concatenate their fractional masses.  If their total mass is $M_i$, slot $z$
on machine $i$ is
\[
 S_{i,z}=[z-1,\min\{z,M_i\}),
 \qquad z=1,\ldots,\lceil M_i\rceil.
\]
A job is adjacent to a slot precisely when their half-open intervals have
nonempty intersection.  A valid matching assigns every job to an adjacent
slot and uses every slot at most once.  The induced schedule assigns the job
to the machine containing its matched slot.

\begin{theorem}\label{thm:spectrum}
For every real $0<\delta<1/6$, this slot graph has exactly eight valid
job-saturating matchings.  Their makespan multiset is
\[
 \left\{\left(\frac{11}{6}-\delta\right)^1,
 \left(\frac43-\delta\right)^2,1^5\right\}.
\]
Consequently, the minimum makespan is $1$ and the maximum makespan is
$11/6-\delta$.
\end{theorem}

\section{Uniform slot graph}

The strict inequalities
\[
 1=p_q>p_h=\frac12>p_r=\frac13-\delta>0
\]
fix the order $q,h,r$ on machine~0.  The fractional masses on machines
$0,1,2$ are, respectively,
\[
 M_0=2+\delta,\qquad M_1=\frac23-\delta,
 \qquad M_2=\frac13.
\]
Write $A=S_{0,1}$, $B=S_{0,2}$, $C=S_{0,3}$,
$D=S_{1,1}$, and $E=S_{2,1}$.  The job and slot intervals are
\[
\begin{array}{c|c|c}
i&\text{job intervals}&\text{slot intervals}\\ \hline
0&q:[0,\frac13+\delta),\quad
  h:[\frac13+\delta,1+\delta),\quad
  r:[1+\delta,2+\delta)
 &A:[0,1),\ B:[1,2),\ C:[2,2+\delta)\\[2pt]
1&q:[0,\frac23-\delta)&D:[0,\frac23-\delta)\\[2pt]
2&h:[0,\frac13)&E:[0,\frac13).
\end{array}
\]

For half-open intervals $[a,b)$ and $[c,e)$, intersection is nonempty if and
only if $b>c$ and $e>a$.  Applying these strict tests gives
\begin{equation}\label{eq:neighborhoods}
 N(q)=\{A,D\},\qquad N(h)=\{A,B,E\},\qquad
 N(r)=\{B,C\}.
\end{equation}
These neighborhoods do not change on the parameter interval.  In particular,
the parameter-sensitive overlaps $hB$ and $rC$ both have length $\delta>0$.
The closest excluded contacts remain excluded: the relevant signed gaps for
$qB$ and $rA$ are $-2/3+\delta<0$ and $-\delta<0$, respectively.  There is
also no slot-count breakpoint because $M_0\in(2,13/6)$ while $M_1,M_2\in(0,1)$.

\section{Exhaustive matching and load calculation}

We now enumerate the injective choices in \eqref{eq:neighborhoods}.  If $q$
uses $A$, then $h$ can use $B$ or $E$.  The first choice forces $r$ to $C$;
the second allows $r$ to use $B$ or $C$.  This gives three matchings.  If $q$
uses $D$, then $h$ can use $A,B$, or $E$, giving respectively two, one, and
two choices for $r$.  These five further matchings exhaust the product of the
three neighborhoods.

The complete calculation is shown in Table~\ref{tab:loads}.  Coordinates in
the first column give the slots of $(q,h,r)$, and $(L_0,L_1,L_2)$ is the
induced machine-load vector.  Different slot matchings are retained even when
they induce the same machine assignment.

\begin{table}[ht]
\centering
\caption{All valid matchings and their exact loads.}
\label{tab:loads}
\begin{tabular}{@{}ccc@{}}
\toprule
Matching & $(L_0,L_1,L_2)$ & Makespan \\ \midrule
$(A,B,C)$ & $(11/6-\delta,0,0)$ & $11/6-\delta$\\
$(A,E,B)$ & $(4/3-\delta,0,1/2)$ & $4/3-\delta$\\
$(A,E,C)$ & $(4/3-\delta,0,1/2)$ & $4/3-\delta$\\
$(D,A,B)$ & $(5/6-\delta,1,0)$ & $1$\\
$(D,A,C)$ & $(5/6-\delta,1,0)$ & $1$\\
$(D,B,C)$ & $(5/6-\delta,1,0)$ & $1$\\
$(D,E,B)$ & $(1/3-\delta,1,1/2)$ & $1$\\
$(D,E,C)$ & $(1/3-\delta,1,1/2)$ & $1$\\ \bottomrule
\end{tabular}
\end{table}

Throughout $0<\delta<1/6$,
\[
 \frac{11}{6}-\delta>\frac43-\delta>1,
\]
and every nonunit load in the final five rows is below $1$.  The table
therefore gives exactly the multiset in Theorem~\ref{thm:spectrum}.  It also
shows the stated maximum.  Independently, the unit-size job $q$ implies that
no schedule can have makespan below $1$, while the final five rows attain that
value.  This completes the proof. \qed

\section{Verification, comparison, and limitations}

The proof is accompanied by a deterministic exact-rational replay.  From the
instance data it reconstructs the job order, fraction intervals, slot counts,
all signed contact tests, the $12$ raw neighborhood tuples, the eight injective
tuples, and their affine load vectors.  It uses rational arithmetic for
expressions of the form $a+b\delta$ and no floating-point tolerance.  Exact
finite certificate designs with a separate reconstruction step are useful in
substantially larger bounded enumerations as well; for example, Zeng et al.
use integer predicates and independently reconstructed finite graphs in their
ternary Borsuk and kissing-code certificates \cite{ZengEtAl2026}.  In the
present three-job calculation, however, the replay is only a reproducibility
check: the symbolic proof above establishes the quantified result.

The scope is deliberately narrow.  Shavit's Theorem~2 already proves that the
first row of Table~\ref{tab:loads} is a valid matching uniformly in $\delta$;
the contribution here is the exhaustive symbolic spectrum and its extrema.
The theorem excludes $\delta=0$ and $\delta=1/6$, depends on the stated
half-open convention and fixed fractional point, and makes no assertion about
other feasible fractional solutions, solver outputs, thresholds, or arbitrary
instances.  In particular, it does not modify the general $11/6$ analysis of
the Wang--Sitters scheme.

\paragraph{Computational and AI assistance.}
OpenAI Codex assisted with primary-literature verification, translation of the
audited argument into manuscript form, exact-replay preparation, and
adversarial consistency checks.  The computational check uses exact rational
arithmetic and is not used as a substitute for the displayed proof.  No proof
assistant was used.

{\small
\bibliographystyle{plain}
\bibliography{references}
}

\end{document}
